\section{An exact individual cubic length at the coefficient prime three}
\label{sec:exact-cubic-descent}

For the fixed cubic character of $269$, the bounded integral
comparison error in AM-H08 can be eliminated.  The key is that the
same ramified inertia which made the local condition automatic in
AC-G22 also detects the global transgression obstruction.

Retain the notation of AC-G22: $K=\mathbb Q(\omega)$,
$E=K(\alpha)$, $\alpha^3=269$, $\chi(g)=g\alpha/\alpha$,
$\mathcal O=\mathbb Z_3[\omega]$, $T=\mathbb Q_3(\omega)$,
$\lambda=\omega-1$, and
$M_e=(T/\mathcal O)(\chi\varepsilon_3^{-e})$.
Let $\Gamma=\operatorname{Gal}(E/K)$, with generator $\sigma$
chosen so that $\chi(\sigma)=\omega$.
The large set $S$ contains the primes above $3$ and $269$.
Strict groups over $K$ with coefficient $M_e$ are the groups
$\mathcal S_e$ of AC-G22.  The untwisted group
\[
 U_e=S_3(E,e)
 =\ker\left(H^1(G_{E,S_3},\mathbb Q_3/\mathbb Z_3(-e))
       \longrightarrow\bigoplus_{w\mid3}
            H^1(E_w,\mathbb Q_3/\mathbb Z_3(-e))\right)
\]
uses the small set $S_3$ consisting only of the primes above $3$.
As in AM-H06, it is identified by inflation with the subgroup
of the large-$S$ cohomology also zero on inertia above $269$.

\begin{theorem}[AC-G26: the exact one-unit descent correction]
\label{AC-G26}
For every positive integer $e$, put
\[
 a_E(e)=\operatorname{length}_{\mathbb Z_3}U_e,\qquad
 \kappa_e=v_3\bigl(\#K_{2e}(\mathcal O_E)\bigr),\qquad
 k_e=1+v_3(e).
\]
Then there is an exact sequence of finite $\mathcal O$-modules
\[
 0\longrightarrow\mathbb F_3\longrightarrow\mathcal S_e
 \xrightarrow{\operatorname{res}}
       (U_e\otimes_{\mathbb Z_3}\mathcal O(\chi))^\Gamma
 \longrightarrow0,                                    \tag{G26.1}
\]
and the target has $\mathcal O$-length $a_E(e)$.
Consequently the finite cyclic group of AC-G22 has exactly
\[
 \mathcal S_e\simeq\mathcal O/\lambda^{m_e},\qquad
 \boxed{m_e=a_E(e)+1=\kappa_e-2k_e+1}.                 \tag{G26.2}
\]
Its residual Kummer-$\omega$ line is precisely the kernel of
restriction in (G26.1).  No RH hypothesis is used.
The theorem determines the comparison exactly; it does not
bound the growth of $\kappa_e$.
\end{theorem}
\begin{proof}
\proofstep{1. The invariant coefficient and the two group-cohomology terms.}
The prime of $K$ above $3$ splits completely in $E/K$,
and $E\cap\mathbb Q(\mu_{3^\infty})=K$, as proved in
AC-G22 and AC-K11.  Hence the cyclotomic image of $G_E$ is
$1+3\mathbb Z_3$.  With $k=k_e$, it follows that
\[
 N_e:=M_e^{G_{E,S}}=M_e[\lambda^{2k}].                  \tag{G26.3}
\]
The action of $\Gamma$ on $N_e$ is exactly $\chi$.
Indeed every cyclotomic value of $G_K$ is in
$1+3\mathbb Z_3$, and its $(-e)$-th power differs from one
by a multiple of $\lambda^{2k}$, so that factor acts trivially
on (G26.3).  Thus $\sigma$ acts by $\omega$, and
$1+\sigma+\sigma^2$ acts by zero.

For a cyclic group of order three, the periodic cochain resolution
alternates $\sigma-1$ and $1+\sigma+\sigma^2$.  Applied here it
gives
\[
 H^1(\Gamma,N_e)=N_e/\lambda N_e\simeq\mathbb F_3,
 \qquad
 H^2(\Gamma,N_e)=N_e[\lambda]\simeq\mathbb F_3.         \tag{G26.4}
\]
These are the terms in the continuous inflation--restriction
five-term sequence for
$1\to G_{E,S}\to G_{K,S}\to\Gamma\to1$.
This is an actual exact sequence with the large set $S$:
$E/K$ is ramified only there, and the maximal extensions
unramified away from $S$ give the indicated subgroup and quotient.

On $G_E$ the character $\chi$ is trivial.  Since $\mathcal O$
is finite free over $\mathbb Z_3$, the restriction coefficient
and its cohomology are the scalar extensions of
$\mathbb Q_3/\mathbb Z_3(-e)$ and its cohomology.  The
$\Gamma$ action on the latter cohomology is its ordinary Galois
action, multiplied by $\chi$ on the coefficient factor.
Consequently the strict invariant subgroup in the middle target
of inflation--restriction is
\[
                 B_e=(U_e\otimes\mathcal O(\chi))^\Gamma.
                                                               \tag{G26.5}
\]

\proofstep{2. Ramified inertia kills the transgression on the strict target.}
Let $I_K$ be absolute inertia at the prime $(269)$ of $K$,
and let $I_E$ be its index-three subgroup for $E$.
Because $E/K$ is totally ramified at this prime, there is a
commutative map of extension sequences
\[
\begin{array}{ccccccccc}
1&\longrightarrow&I_E&\longrightarrow&I_K&\longrightarrow&\Gamma&\longrightarrow&1\\
 &&\downarrow&&\downarrow&&\Vert\\
1&\longrightarrow&G_{E,S}&\longrightarrow&G_{K,S}&\longrightarrow&\Gamma&\longrightarrow&1.
\end{array}                                             \tag{G26.6}
\]
The vertical maps mean the usual local-to-global homomorphisms;
they need not be injective to use naturality of cohomology.
On $I_E$ both $\chi$ and $\varepsilon_3$ are trivial, so
$M_e^{I_E}=M_e$.  The quotient $\Gamma$ acts on this group
by $\omega$, because the cyclotomic character is unramified
at $269$.  Therefore
\[
 H^2(\Gamma,M_e)=M_e[\lambda],
 \qquad
 H^2(\Gamma,N_e)\longrightarrow H^2(\Gamma,M_e)
 \quad\hbox{is an isomorphism}.                         \tag{G26.7}
\]
The last assertion follows directly from (G26.4): it is the
inclusion of the same one-dimensional $\lambda$-socle, and
both norm operators are zero.

Let $b\in B_e$.  Its restriction to $H^1(I_E,M_e)$ is zero,
since $U_e$ is unramified outside $3$.  Naturality of the
five-term sequence for (G26.6) sends the global transgression
$d_2(b)\in H^2(\Gamma,N_e)$ to the local transgression of
zero.  By the injectivity in (G26.7), $d_2(b)=0$.
Thus every $b\in B_e$ lifts to a class of
$H^1(G_{K,S},M_e)$.

Every such lift is strict.  At a prime above $3$, restriction
from $K$ to any chosen completion of $E$ is restriction along
an isomorphism $E_w=K_3$, since $3$ splits completely; the
coefficient is the same there because $\chi$ is locally trivial.
As $b$ is locally zero, so is its lift.  At $269$, AC-G22
proves $H^1(I_K,M_e)=0$ using the nontrivial tame coefficient
action and the surjectivity of $\lambda$ on $T/\mathcal O$.
This verifies the remaining strict condition.

Conversely a strict class restricts into $B_e$: at $3$ this is
immediate, and at $269$ the restriction of a local coboundary
to $I_E$ is zero because $I_E$ acts trivially.  The resulting
large-$S$ class therefore descends from the small-$S_3$ group
over $E$.  Finally, the inflated subgroup
$H^1(\Gamma,N_e)$ is strict.  Its restrictions at $3$ are zero
because the local quotient to $\Gamma$ is trivial, and its
inertia condition at $269$ is automatic as above.
Inflation--restriction and (G26.4) now prove (G26.1).
Its kernel is an $\mathcal O$-line of length one, and by AC-G22
that line must be the unique socle, namely the Kummer-$\omega$
line.

\proofstep{3. The target has exactly the untwisted length.}
Corestriction takes $U_e$ to the untwisted strict group over
$K$, which is zero by the AC-G22 base corollary.  To verify the
claim about the local conditions despite the ramification of
$E/K$, first work in the large-$S$ cohomology groups.
At $3$, the Mackey formula expresses localization of
corestriction as the sum of the local corestrictions of the zero
classes.  At $269$, apply the same formula to inertia:
restriction to $I_E$ is zero, so its corestriction to $I_K$
is zero.  The coefficient is unramified at $269$.
Inflation--restriction then identifies the corestricted class
with a class from $G_{K,S_3}$, still zero at $3$.
It is therefore in the zero strict group over $K$.

Restriction of corestriction is the group norm, so
\[
           (1+\sigma+\sigma^2)U_e=0.                  \tag{G26.8}
\]
It follows that $U_e$ is a module over
\[
 R=\mathbb Z_3[\Gamma]/(1+\sigma+\sigma^2)
        \simeq\mathbb Z_3[\omega].
\]
This statement does not require an integral character projector.
Put $P=U_e\otimes_{\mathbb Z_3}\mathcal O$ and let
$q=\sigma\otimes\omega-1$, an $\mathcal O$-linear endomorphism
of the finite module $P$.  Its kernel is $B_e$.
Its cokernel is
\[
 P/qP\simeq U_e\otimes_{R,\,\sigma\mapsto\omega^{-1}}
                        \mathcal O.                    \tag{G26.9}
\]
Indeed the quotient imposes precisely the relation
$\sigma u\otimes a=u\otimes\omega^{-1}a$.
The ring map in (G26.9) is an isomorphism; it only conjugates
the two cubic roots.  Hence the cokernel has the same cardinality
as $U_e$.  An endomorphism of a finite group has kernel and
cokernel of equal cardinality, so
\[
 \#B_e=\#U_e=3^{a_E(e)},\qquad
 \operatorname{length}_{\mathcal O}B_e=a_E(e).          \tag{G26.10}
\]
All finiteness statements here follow from AC-G22 or AC-K10;
no growth assertion is hidden in taking these cardinalities.

\proofstep{4. The exact numerical identity.}
Taking $\mathcal O$-lengths in (G26.1) gives
$m_e=1+a_E(e)$.  The exact Poitou--Tate/K-theory formula
AC-K11 gives $a_E(e)=\kappa_e-2k_e$.  Substitution proves
(G26.2).
\end{proof}

\begin{corollary}[Exact finite-weight correction]
\label{AC-G26-sum}
For every integer $t\geq2$, writing $N=t-1$,
\[
 \sum_{1\leq e<t}m_e
 =\sum_{1\leq e<t}\kappa_e-N-2v_3(N!)
 =\sum_{1\leq e<t}\kappa_e-2N+s_3(N),                 \tag{G26.11}
\]
where $s_3(N)$ is the sum of the base-three digits of $N$.
Multiplication by $\log3$ gives the exact prime-$3$ contribution
in the logarithmically weighted cohomology budget.
\end{corollary}
\begin{proof}
Sum (G26.2) and use
$\sum_{e=1}^Nv_3(e)=v_3(N!)$.
Counting the multiples of $3^j$ gives
$v_3(N!)=\sum_{j\geq1}\lfloor N/3^j\rfloor
=(N-s_3(N))/2$, proving the second equality.
\end{proof}

\begin{remark}[What has been reduced]
The coefficient-prime-$3$ comparison now has no unspecified
constant and no Fourier-lattice defect: the additional length is
exactly one.  Estimating $m_e$ is consequently equivalent, with
the explicit local subtraction in (G26.2), to estimating
$v_3\#K_{2e}(\mathcal O_E)$ for the one fixed degree-six field.
This is a smaller concrete arithmetic problem than bounding all
character components of an arbitrary host.  It is not an estimate
of those K-group orders, and it does not itself replace the
varying fields in \texttt{VWinAt} or prove the all-window
inequality (H04.4).
\end{remark}
